\frontchapter{Предисловие}

Контейнеры изменили то, как программы собирают, тестируют и запускают.
Сегодня \enquote{упаковать приложение в Docker}~--- такой же базовый навык разработчика
и системного администратора, как работа с Git или командной строкой.
Эта книга проведёт вас от первого \cmd{docker run} до осознанного
проектирования образов, сетей и конфигурации \cmd{docker compose},
попутно объясняя, что именно происходит внутри.

\section*{Кому адресована книга}

Книга рассчитана на читателя, который умеет работать в терминале и имеет
базовое представление о Linux: файловая система, процессы, права доступа,
переменные окружения. Знание языков программирования не требуется~---
примеры намеренно простые (статический сайт, небольшой сервис на Python
или Node.js), а основное внимание уделено Docker.

Материал будет полезен:

\begin{itemize}
  \item разработчикам, которые хотят, чтобы приложение одинаково работало
        на ноутбуке, в тестовой среде и в production;
  \item администраторам и DevOps-инженерам, которым нужны воспроизводимые
        сборки, понятная диагностика и безопасные конфигурации;
  \item студентам и всем, кто готовится к собеседованиям, где спрашивают
        про слои образа, \cmd{ENTRYPOINT}, тома и сети;
  \item преподавателям: готовые лабораторные работы с пошаговыми решениями
        можно использовать как практикум.
\end{itemize}

\section*{Что нужно для работы}

\begin{itemize}
  \item Компьютер под Linux, macOS или Windows с установленным Docker
        (установка описана в главе~3). Для Windows и macOS проще всего
        начать с Docker Desktop; для Linux~--- с Docker Engine.
  \item Терминал (bash, zsh или PowerShell) и текстовый редактор.
  \item Желательно~--- учётная запись на Docker Hub: она понадобится
        в главе~10, но не раньше.
  \item Достаточно 8~ГБ оперативной памяти и 10--15~ГБ свободного места
        на диске.
\end{itemize}

Все команды книги проверены на Docker Engine~25 и новее и Docker Compose~v2.
Отдельные флаги могли появиться в более поздних версиях~--- если команда
не распознаётся, сравните вывод \cmd{docker version} с документацией.
Контейнерные интерфейсы развиваются медленно, поэтому основной материал
останется актуальным надолго.

\section*{Как устроена книга}

Книга состоит из четырёх частей. Первые три читаются последовательно,
четвёртую можно использовать как справочник.

\begin{center}
\small
\begin{tabularx}{\textwidth}{@{}lX@{}}
\toprule
\textbf{Часть} & \textbf{Содержание} \\
\midrule
I. Введение & Зачем нужны контейнеры, как устроена изоляция, установка
              и первый запуск. \\
II. Образы и контейнеры & Слои, теги, жизненный цикл, Dockerfile,
              постоянные данные. Ядро книги. \\
III. Сети и оркестрация & Драйверы сетей, публикация портов, Docker Compose
              для многоконтейнерных приложений. \\
IV. Эксплуатация & Реестры и CI/CD, диагностика, безопасность,
              оптимизация образов. \\
\multicolumn{2}{@{}l@{}}{\itshape Планируется: практикум, шпаргалка команд, глоссарий.} \\
\bottomrule
\end{tabularx}
\end{center}

Если вы совсем новичок, идите по главам подряд и повторяйте примеры
после каждой части. Если вы уже пользуетесь Docker, начните с глав~6, 7 и 12:
там собраны типичные ошибки и правила, которые чаще всего нарушают.

\section*{Обозначения}

В книге используются пять типов выделенных блоков.

\begin{defbox}
Важное понятие или точная формулировка, на которую можно ссылаться.
\end{defbox}

\begin{notebox}
Пояснение, историческая справка или деталь, которую полезно знать,
но можно пропустить при первом чтении.
\end{notebox}

\begin{warnbox}
Распространённая ошибка, источник проблем с безопасностью или данными.
\end{warnbox}

\begin{tipbox}
Практический приём, экономящий время.
\end{tipbox}

\begin{labbox}{Пример оформления}
Так будут выглядеть лабораторные работы практикума: цель, шаги,
проверка результата и разбор.
\end{labbox}

\section*{Соглашения о командах}

Команды, которые нужно ввести в терминал, начинаются с приглашения
\cmd{\$}. Строки без приглашения~--- это вывод программы или содержимое
файла, их вводить не нужно.

\begin{code}[language=bash]
$ docker run --rm hello-world
Unable to find image 'hello-world:latest' locally
...
Hello from Docker!
\end{code}

Многоточие \cmd{...} внутри вывода означает пропущенный фрагмент.
Имена вроде \cmd{myapp}, \cmd{web} или \cmd{8080}~--- примеры; заменяйте их
на свои. В листингах Dockerfile и YAML комментарии начинаются с~\cmd{\#}.

Если команда выполняется от имени обычного пользователя, а Docker установлен
так, что требует \cmd{sudo}, вы увидите ошибку доступа к сокету. Как это
исправить~--- в разделе~3.4.

\section*{Как заниматься}

Лучший способ освоить Docker~--- набирать команды руками, а не копировать.
После каждой главы есть блок «Итоги» и вопросы для самопроверки: если
ответы не складываются, вернитесь к соответствующему разделу.
Все примеры стоит повторить руками: от одиночного контейнера
до приложения из нескольких сервисов с базой данных, томами
и проверками здоровья.

\begin{tipbox}
Заведите отдельный каталог для упражнений, например \file{\textasciitilde/docker-lab},
и держите его под Git. Так вы увидите, как меняются Dockerfile и compose-файлы,
и сможете вернуться к удачному варианту.
\end{tipbox}

Желаю интересной работы. Контейнеры устроены проще, чем кажутся,
если разобраться в нескольких базовых идеях~--- именно им посвящены
первые главы.
